\documentclass[10pt,a4paper]{amsart}
\usepackage[margin=19mm]{geometry}
\usepackage{amsmath,amssymb,mathtools}
\usepackage[scr=rsfs,cal=euler]{mathalfa}
\usepackage{xfrac}
\usepackage[unicode,hidelinks,bookmarks=false]{hyperref}
\newcommand{\ie}{i.e.\ }
\newcommand{\cf}{cf.\ }
\newcommand{\resp}{respectively\ }
\newcommand{\Subsection}{\subsection}
\providecommand{\nobreakdash}{\nobreak\hbox{-}}
\setlength{\emergencystretch}{2em}
\multlinegap0pt
\usepackage[normalem]{ulem}
\usepackage{amssymb}
\usepackage{mathtools}
\usepackage{braket}
\usepackage{bm}
\usepackage{enumerate}
\usepackage{xcolor}
\usepackage{cancel}
\usepackage[shortlabels]{enumitem}
\newtheorem{prop}{Proposition}
\newtheorem{defn}{Definition}
\newcommand{\be}{\begin{equation}}
\newcommand{\ee}{\end{equation}}
\title{Convergence of Timed-Metric Spaces and Causality}
\author{Mauricio Che, Raquel Perales}
\date{Source: arXiv:2604.22362v1 (CC BY 4.0)}
\begin{document}
\maketitle

\noindent\emph{Source and licence.} Convergence of Timed-Metric Spaces and Causality. Version arXiv:2604.22362v1 (CC BY 4.0), distributed by arXiv under the Creative Commons Attribution 4.0 International licence, http://creativecommons.org/licenses/by/4.0/, as recorded in the arXiv licence metadata for this exact version and shown on the abstract page https://arxiv.org/abs/2604.22362v1. Full licence notice: \texttt{licenses/arxiv-2604.22362-CC-BY-4.0.txt}.\par\smallskip

\noindent\textit{Editorial scope.} Counted unit: the article's prop greater (source0.tex lines 529--537). The article's complete original proof of it is delivered with it (source0.tex lines 1804--1862, one \texttt{proof} environment of 59 lines, which the article places away from the statement). No mathematics is written, repaired or re-derived: the statement and the proof are the article's own lines, in the article's own notation, and no retained line is substituted or re-flowed.  Retained uncounted as the source-local context the proof needs: lines 1770--1787.  Nothing was omitted from any retained span, so the package carries no omission record.  No retained line contains a citation, so the wrapper carries no bibliography and no bibliography entry.  No retained line contains a figure environment, an image-inclusion command or drawing code, so no image is delivered and the package declares no asset dependency.  Editorial formatting only, in addition: an AMS \texttt{amsart} preamble that restates the article's own declarations and macros used by the retained lines, namely the environments \texttt{prop}, \texttt{defn} on separate counters, together with the standard packages those lines need.  The retained environments print in this document as Proposition 1, Definition 1 in order of appearance, whereas the article prints the same items with the numbering of its own full document.  The complete native source of the article as distributed in the arXiv e-print for this exact version is bundled unchanged as \texttt{sources/199083/source0.tex}.
\begin{defn}\label{defn:null-distance}
Given a timed-metric space $(X,d,\tau)$, 
we define $\hat{d}_{d,\tau}: X \times X \to [0, \infty]$
as
\be\label{eq:null}
\hat{d}_{d,\tau}(p,q)=\inf \sum_{i=1}^{N} |\tau(p_i)-\tau(p_{i-1})| 
\ee
where the infimum is over all collections of points, if they exist, $\{p_i\}_{i=0}^{N}$, where  
$p=p_{0}$ and $p_{N}=q$, and $p_i$, $p_{i-1}$ are causally related, possibly alternating past
with future:
\be\label{eq:null-points}
p_i\in J^-(p_{i-1}) \quad \text{ or } \quad p_i \in J^+(p_{i-1}), \quad 
 i=1,\ldots,N.
\ee
Otherwise, $\hat{d}_{d,\tau}(p,q)=\infty$. 
We call $\hat{d}_{d,\tau}$ the  associated \textbf{null distance} of 
$(X,d,\tau)$, although it is an extended distance in general. If $\hat{d}_{d,\tau}=d$ we say that $(X,d,\tau)$ is \textbf{causally null}.  
\end{defn}

\begin{prop}\label{prop-null-distance is greater}
Let $(X,d,\tau)$ be a timed-metric space. 
Then the null distance $\hat{d}_{d,\tau}$, given in Definition~\ref{defn:null-distance}, is an extended distance such that
\begin{equation}\label{eq:null-distance is greater}
\hat{d}_{d,\tau}\geq d,
\end{equation}
and $\tau$ is $1$-Lipschitz with respect to $\hat{d}_{d,\tau}$.
Moreover, if $\hat{d}_{d,\tau}$ only takes finite values then $(X,\hat{d}_{d,\tau},\tau)$ is a timed-metric space and $\hat{d}_{\hat{d}_{d,\tau},\tau} = \hat{d}_{d,\tau}$.
\end{prop}

\begin{proof}[Proof of Proposition \ref{prop-null-distance is greater}]

In what follows, to avoid confusion, let us denote  $\hat{d}_{d,\tau}$ as $\hat{d}$, and the causal structure of $(X,\hat d,\tau)$ as $J^\pm_{\hat d}$. 

First we prove that $\hat{d}$ is an extended distance. It is clear that $\hat{d}$ is non-negative and symmetric. Let $x,y,z\in X$ and without loss of generality assume that 
$\hat{d}(x,y), \hat{d}(y,z)< \infty$.
Then for any $\varepsilon>0$ let $\{p_i\}_{i=0}^{N}$ and $\{q_i\}_{i=0}^{M}$ be such that $p_0=x$, $p_N = y = q_0$, $q_M=z$, and $p_i\in J^\pm(p_{i-1})\cap J^\pm(p_{i+1})$ and $q_i\in J^\pm(q_{i-1})\cap J^\pm(q_{i+1})$, and 
\be
\sum_{i=1}^{N} |\tau(p_i)-\tau(p_{i-1})| < \hat{d}(x,y) + \varepsilon/2 
\ee
and
\be
\sum_{i=1}^{M} |\tau(q_i)-\tau(q_{i-1})| < \hat{d}(y,z) + \varepsilon/2.
\ee
Then
\be
\hat{d}(x,z) \leq \sum_{i=1}^{N} |\tau(p_i)-\tau(p_{i-1})| +  \sum_{i=1}^{M} |\tau(q_i)-\tau(q_{i-1})| \leq \hat{d}(x,y) + \hat{d}(y,z) + \varepsilon.
\ee
By letting $\varepsilon\to 0$, it follows that $\hat d$ satisfies the triangle inequality. 


Now take $x,y \in X$. By definition,
\begin{align*}
\hat{d}(x,y) = \inf\left\{\sum_{i=1}^{N} |\tau(p_i)-\tau(p_{i-1})|: x = p_0,\dots, p_{N} = y,\ p_{i}\in J^{\pm}(p_{i-1})\right\}.
\end{align*}
Since $p\in J^{+/-}(q)$ is equivalent to $|\tau(q)-\tau(p)| = d(q,p)$, we get
\be
\sum_{i=1}^{N} |\tau(p_i)-\tau(p_{i-1})| = \sum_{i=1}^{N} d(p_i,p_{i-1}) \geq d(x,y)
\ee
whenever $x = p_0,\dots, p_{N} = y$, $p_{i}\in J^\pm_{d,\tau}(p_{i-1})$. Taking the infimum over such tuples $\{p_i\}_{i=0}^{N}$, inequality \eqref{eq:null-distance is greater} follows. In particular, \eqref{eq:null-distance is greater} implies that $\hat{d}$ is positive definite, and 
that $\tau$ is $1$-Lipschitz with respect to $\hat{d}$, given that $\tau$ is $1$-Lipschitz with respect to $d$.




Now assume that $\hat{d}(x,y)<\infty$ for any $x,y\in X$. Therefore $(X,\hat{d},\tau)$ is a timed-metric space. To prove that $\hat{d}_{\hat{d},\tau}= \hat d$ in this case, we first show that the causal structures agree:
\begin{equation}\label{eq-structuresAgree}
    J^-(y)=J^-_{\hat{d}}(y) \quad \forall y \in X. 
\end{equation}
Indeed, let $x\in J^-_{d}(y)$. Taking $p_0 = p_1 = x$ and $p_2 = y$ in the definition of $\hat{d}(x,y)$ yields 
\begin{align*}
\hat{d}(x,y) \leq & (\tau(x)-\tau(x)) + (\tau(y) - \tau(x)) \\
= & d(x,y) \leq \hat{d}(x,y), 
\end{align*}
where in the last part we used \eqref{eq:null-distance is greater}.
Hence, all inequalities are equalities. In particular,  $\hat{d}(x,y) = \tau(y)-\tau(x)$, which proves $x\in J^-_{\hat{d}}(y)$.  Conversely, if $x\in J^-_{\hat{d}}(y)$, by \eqref{eq:null-distance is greater} and $\tau$  $1$-Lipschitz with respect to $d$, we get
\begin{align*}
\tau(y) - \tau(x) = & \hat{d}(x,y)  
\geq  d(x,y) \geq \tau(y)-\tau(x).
\end{align*}
Therefore, $d(x,y) = \tau(y)-\tau(x)$, i.e., 
$x\in J^-(y)$. This concludes the proof of \eqref{eq-structuresAgree}. Thus, 
\begin{align*}
\hat{d}_{\hat{d}, \tau}(x,y) 
&= \inf\left\{\sum_{i=1}^{N} |\tau(p_i)-\tau(p_{i-1})| : x=p_0,\dots, p_{N}=y,\ p_{i}\in J^\pm_{\hat{d}}(p_{i-1})\right\}\\
&= \inf\left\{\sum_{i=1}^{N} |\tau(p_i)-\tau(p_{i-1})| : x=p_0,\dots, p_{N}=y,\ p_{i}\in J^\pm(p_{i-1})\right\}\\
&= \hat{d}(x,y).
\end{align*}
\end{proof}

\end{document}
